\documentclass{article}

\usepackage{fancyhdr}
\usepackage[T1]{fontenc}
\usepackage[margin=0.2in]{geometry}
\usepackage{parskip}
\usepackage{amsmath}
\usepackage{amssymb}

\pagestyle{fancy}
\fancyhead[l]{Andrew O'Dwyer}
\fancyhead[c]{Ritangle Question 17}
\fancyhead[r]{\today}

\begin{document}

\section{Problem Description}
\huge
p, q, r are random integer between 0, 6 inclusive.

these values represent the ratio of angle of a triangle
inside a circle with its verticies on the circles 
circumferance.

what is the probability that the centre of the circle 
isn't in the triangle/on the triangle edge.

\section{Problem Solution}
\[k (p + q + r) = 180\]

The largest angle must be greater than 90 degrees following 
something like the circle theorem that states the angle is
90 degrees when it line passes through the centre of the 
circle.

This means that if p is the largest angle, then:
\begin{align*} 
    \frac{p}{p + q + r} > 0.5 \\
    \implies p > (q+r)
\end{align*} 

The dice can be moddeled as a discrete random distribution.
So the PGF would be:
\begin{align*}
    G_p(t)=G_q(t)=G_r(t)\\
    = \frac{1}{6}(t^6+t^5+t^4+t^3+t^2+t)\\
\end{align*}

And then q+r,p+q and p+r can be moddeled by the PGF:
\begin{align*}
    &G_q(t) \times G_r(t) = G_q(t) \times G_p(t)
        = G_p(t) \times G_r(t)\\
    &= \left( \frac{1}{6}(t^6+t^5+t^4+t^3+t^2+t) \right)^2\\
    &= \frac{1}{36}\left(t^6+t^5+t^4+t^3+t^2+t\right)^2\\
    &= \frac{1}{36}\left(t^{6}+t^5+t^4+t^3+t^2+t\right)
        \left(t^6+t^5+t^4+t^3+t^2+t\right)\\
    &=\frac{1}{36}(t^{12}+t^{11}+t^{10}+t^9+t^8+t^7)\\
    &+\frac{1}{36}(t^{11}+t^{10}+t^9+t^8+t^7+t^6)\\
    &+\frac{1}{36}(t^{10}+t^9+t^8+t^7+t^6+t^5)\\
    &+\frac{1}{36}(t^9+t^8+t^7+t^6+t^5+t^4)\\
    &+\frac{1}{36}(t^8+t^7+t^6+t^5+t^4+t^3)\\
    &+\frac{1}{36}(t^7+t^6+t^5+t^4+t^3+t^2)\\
    &=\frac{1}{36}(t^{12}+2t^{11}+3t^{10}+4t^9+5t^8+6t^7\\
    &+5t^6+4t^5+3t^4+2t^3+t^2)\\
\end{align*}

So now we just need to find the probability that the largest
angle is greater than the sum of the other 2 angles. To do
this I can just check the cases where one of the angles is
greater than the other 2 (as they are all same distributions).
So assuming p is the largerst angle ratio, I should have:
\[ 
    P(q+r=i) = 
    \begin{cases}
        \frac{i-1}{36}  & 1<=i<=6\\
        \frac{13-i}{36} & 7<=i<=12\\
    \end{cases} 
\]

\[ 
    P(p>i) = 
    \begin{cases}
        \frac{6-i}{6} & 1<=i<=6\\
        0             & 7<=i<=12\\
    \end{cases} 
\]

\[ 
    P(p>q+r) = \sum_{i=1}^{12} (P(q+r=i) \times P(p>i))\\
\]

\begin{align*}
    P(p>q+r) &= \sum_{i=1}^{12} (P(q+r=i) \times P(p>i))\\
    &= \sum_{i=1}^{6} 
        \left( \frac{i-1}{36}\times \frac{6-i}{6} \right)
    + \sum_{i=7}^{12}
        \left( \frac{13-i}{36}\times 0 \right)\\
    &= \sum_{i=1}^{6} 
        \left( \frac{(i-1)(6-i)}{216}\right)\\
    &= \sum_{i=1}^{6} 
        \left( \frac{-i^2+7i-6}{216})\right)
\end{align*}

\text{Now decompose and evaluate summations:}

\begin{align*}
    &= \frac{- 1}{216}\sum_{i=1}^{ 6} (i^2)
     + \frac{  7}{216}\sum_{i=1}^{ 6} (i)
     + \frac{- 6}{216}\sum_{i=1}^{ 6} (1)\\
    %
    &= \frac{- 1\times 6}{216\times6}(6+1)(12+1)\\
    &+ \frac{  7\times 6}{216\times2}(6+1)\\
    &+ \frac{- 6\times (6-1)}{216}\\
    %
    &= \frac{13}{108}
\end{align*}

\end{document}
